% !TEX root = ../Thesis.tex
% ============================================================
\chapter{G-Retriever and Prize-Collecting Steiner Trees} \label{chapter:pcst}
% ============================================================

The behaviour this thesis observes on clinical data follows from G-Retriever's
prize and cost definitions rather than from anything clinical, and the only way
to show that is to write the definitions down. Chapter~\ref{chapter:background}
placed PCST-based selection among the alternatives without the algebra. The
algebra follows here, and Section~\ref{sec:pcst-consequences} draws predictions
from it that Chapter~\ref{chapter:results} then puts to the test.

\section{Subgraph Selection as an Optimisation Problem} \label{sec:selection-formal}

Let $G = (V, E)$ be a graph in which every node $v \in V$ and every edge
$e \in E$ carries a text description. A sentence encoder~\cite{reimers2019sentencebert}
maps each description to a vector, giving $\mathbf{x}_v$ for nodes,
$\mathbf{x}_e$ for edges, and $\mathbf{q}$ for the question. Relevance is
measured by cosine similarity,
%
\begin{equation}
\operatorname{sim}(\mathbf{a}, \mathbf{b})
  = \frac{\mathbf{a}^{\top}\mathbf{b}}{\lVert \mathbf{a}\rVert \, \lVert \mathbf{b}\rVert},
\label{eq:cosine}
\end{equation}
%
which is the only place the question enters the computation. Everything after
this point operates on the resulting scores, a fact worth holding onto: no stage
of the algorithm can recover information that Equation~\ref{eq:cosine} failed to
capture, and Section~\ref{sec:pcst-consequences} returns to what that costs when
many nodes share identical text.

A selection strategy returns a subgraph $S = (V_S, E_S)$ with $V_S \subseteq V$
and $E_S \subseteq E$. The two requirements in tension are that $S$ should
contain the entities an answer needs, and that $S$ should be small. Stating them
as one objective, rather than as a filter followed by a repair step, is what
distinguishes the PCST formulation from the families surveyed in
Section~\ref{sec:subgraph-selection-approaches}.

\section{The Prize-Collecting Steiner Tree} \label{sec:pcst-definition}

The Steiner tree problem asks for a minimum-cost tree spanning a set of required
terminals. Its prize-collecting variant~\cite{bienstock1993pcst} relaxes the
requirement: terminals may be left out, at the price of forfeiting a prize
attached to them. Formally, given prizes $p: V \to \mathbb{R}_{\geq 0}$ and costs
$c: E \to \mathbb{R}_{\geq 0}$, the problem is to find a connected subgraph $S$
maximising
%
\begin{equation}
F(S) \;=\; \sum_{v \in V_S} p(v) \;-\; \sum_{e \in E_S} c(e).
\label{eq:pcst-objective}
\end{equation}
%
Equivalently, and in the form the classical literature states it, $S$ minimises
the forfeited prize plus the cost paid,
$\sum_{v \notin V_S} p(v) + \sum_{e \in E_S} c(e)$; the two differ by the
constant $\sum_{v \in V} p(v)$.

Three features of Equation~\ref{eq:pcst-objective} matter for retrieval. First,
connectivity is a constraint on the feasible set, not a term in the objective, so
it cannot be traded away by a sufficiently attractive node. Second, the trade-off
between relevance and size is priced by a single quantity, the cost per edge, and
the size of the result is therefore controlled by a continuous parameter rather
than by an integer count. Third, and least obvious, the formulation admits nodes
of low prize: a node with $p(v) = 0$ is worth including whenever it lies on a path
whose endpoints jointly forfeit more than the connecting edges cost. Such a node
is invisible to any method that scores nodes independently, and it is the
mechanism by which a bridging entity can reach the prompt at all.

\section{G-Retriever's Prize and Cost Assignment} \label{sec:gretriever-prizes}

Equation~\ref{eq:pcst-objective} says nothing about where prizes come from.
G-Retriever~\cite{he2024gretriever} supplies them by rank rather than by score.

\paragraph{Node prizes.} Let $\pi(v)$ be the rank of node $v$ when nodes are
sorted by $\operatorname{sim}(\mathbf{q}, \mathbf{x}_v)$ in descending order,
with $\pi(v) = 1$ for the most similar. For a parameter $k$ (\texttt{topk}),
%
\begin{equation}
p(v) \;=\;
\begin{cases}
k - \pi(v) + 1, & \pi(v) \leq k, \\
0, & \text{otherwise.}
\end{cases}
\label{eq:node-prizes}
\end{equation}
%
So the best node receives $k$, the next $k-1$, down to $1$ for the $k$-th, and
every other node receives nothing. Using rank rather than the similarity value
makes the prize scale independent of how tightly the similarities are clustered,
which is a reasonable defence against an encoder whose scores occupy a narrow
band. It also discards magnitude entirely: the gap between the best and second
best node is always exactly one prize unit, whether the encoder considered them
near-identical or wildly different.

\paragraph{Edge prizes.} Edges are ranked the same way, by
$\operatorname{sim}(\mathbf{q}, \mathbf{x}_e)$, with a parameter $k_e$
(\texttt{topk\_e}) playing the role of $k$. The difference is what happens to
ties. Let $\tau$ index the distinct similarity values among the top $k_e$, let
$b(\tau)$ be the prize budget assigned to tier $\tau$ by the same descending
rule, and let $n_\tau$ be the number of edges whose similarity equals that
value. Every edge in tier $\tau$ receives
%
\begin{equation}
p(e) \;=\; \frac{b(\tau)}{n_\tau}.
\label{eq:edge-prizes}
\end{equation}
%
Dividing the tier's budget among tied edges is a sensible way to stop a single
relation type from dominating the objective simply because it is common. Its
consequences on a graph where relation types \emph{are} overwhelmingly common are
the subject of Section~\ref{sec:pcst-consequences}.

\paragraph{Edge cost.} Every edge is charged the same constant $c_e$
(\texttt{cost\_e}), so
%
\begin{equation}
\sum_{e \in E_S} c(e) \;=\; c_e \, \lvert E_S \rvert ,
\label{eq:edge-cost}
\end{equation}
%
and the cost term is proportional to the number of edges retrieved. This is the
size dial: raising $c_e$ makes each connection more expensive and shrinks the
result.

\section{Making Edges Carry Prizes} \label{sec:virtual-nodes}

Equation~\ref{eq:pcst-objective} places prizes on nodes and costs on edges, but
Equation~\ref{eq:edge-prizes} places prizes on edges as well. Solvers for the
prize-collecting Steiner tree do not accept that input, so G-Retriever removes
the difficulty by transformation rather than by changing the solver.

Consider an edge $e = (u, w)$ with prize $p(e)$ and cost $c_e$. Two cases arise.

If $p(e) \leq c_e$, the prize is absorbed into the cost. The edge is kept as a
single edge with adjusted cost
%
\begin{equation}
c'(e) \;=\; c_e - p(e) \;\geq\; 0 ,
\label{eq:absorbed-cost}
\end{equation}
%
which is exact: selecting the edge gains $p(e)$ and pays $c_e$, a net payment of
$c_e - p(e)$, and the objective is unchanged by writing it this way.

If $p(e) > c_e$, absorbing the prize would require a negative cost, which the
solver cannot represent. Instead a \emph{virtual node} $\nu_e$ is introduced,
carrying prize $p(e) - c_e$, and the edge is replaced by two edges
$(u, \nu_e)$ and $(\nu_e, w)$, each of cost $c_e / 2$. Selecting both new edges
costs $c_e$ in total and collects $p(e) - c_e$, giving a net gain of
$p(e) - c_e$, which again matches the original. The transformation is therefore
objective-preserving, and it costs one additional node and one additional edge
per high-prize edge.

The price of the construction is paid at decoding time. The solver returns a set
of vertices in the transformed graph, some real and some virtual, and the virtual
ones must be mapped back to the edges they stand for. The final edge set is the
union of the real edges the solver chose and the edges recovered from selected
virtual nodes. That union is assembled by concatenation and is not re-sorted,
which is harmless for a set-based metric and misleading for anything comparing
edge sequences positionally; this thesis normalises the ordering at the interface
rather than in the algorithm, so that the implementation stays faithful to the
published one (Section~\ref{sec:implementation-fidelity}).

\section{Keeping Costs Below the Largest Prize} \label{sec:cost-adjustment}

One further adjustment sits between the parameters and the solver. If the
per-edge cost exceeds every edge prize, the transformation of
Section~\ref{sec:virtual-nodes} never takes its second branch, and the edge
prizes cease to influence the solution at all. To avoid that degenerate regime,
the implementation caps the cost at
%
\begin{equation}
c_e \;\leftarrow\; \min\!\left\{\, c_e, \;\; \left(\max_{e \in E} p(e)\right)\left(1 - \frac{\gamma}{2}\right) \right\},
\qquad \gamma = 0.01 ,
\label{eq:cost-adjustment}
\end{equation}
%
so that the cost is held just below the largest edge prize in the graph.

Equation~\ref{eq:cost-adjustment} is easy to read as a safety rail and hard to
notice as a coupling. It ties the size dial to the edge prizes: whatever the user
requests, the effective cost can never exceed $\max_e p(e)$. Two consequences
follow immediately, and both are tested in Chapter~6. Raising $c_e$ above that
maximum has no effect whatsoever, so the dial saturates from above. And if
$\max_e p(e)$ is itself very small, the effective cost is dragged down with it,
so the dial collapses regardless of what value was requested.

\section{Solving and Decoding} \label{sec:pcst-solver}

The prize-collecting Steiner tree problem is NP-hard, and G-Retriever uses the
primal-dual approximation of Goemans and Williamson in the near-linear-time
formulation of Hegde et al.~\cite{hegde2015nearlylinear}, which is what the
\texttt{pcst\_fast} package implements. The solver is invoked unrooted, asked for
a single connected component, and run with Goemans-Williamson pruning. Unrooted
is the right choice here because no node is distinguished a priori: in a rooted
formulation the tree must contain a designated node, which would amount to
assuming the answer's location.

Decoding proceeds in three steps: partition the returned vertices into real and
virtual, map virtual vertices back to their edges, and take the union of that set
with the real edges the solver returned. The resulting node and edge sets index
the original graph, and the retained rows of the node and edge tables are
serialised into the text that reaches the language model.

\section{What the Definitions Predict on a Clinical Graph} \label{sec:pcst-consequences}

The definitions above were designed for graphs with many relation types used a
few times each and with text that distinguishes one node from another. The graph
measured in Section~\ref{sec:problem-statement} has neither property. Three
predictions follow from the algebra alone, before any experiment.

\paragraph{Prediction 1: the edge-prize mechanism dilutes to nothing.} Combining
Equations~\ref{eq:edge-prizes} and~\ref{eq:cost-adjustment} gives the effective
cost as a function of the largest tier. Since text is shared across every use of
a relation type, all $n_\tau$ edges of a type receive identical similarity and
therefore fall in one tier. In the measured graph the relation \texttt{hasCode}
is used $4{,}883{,}809$ times, so with a top tier budget of $b \le k_e$ the
per-edge prize is at most
%
\begin{equation}
p(e) \;=\; \frac{b}{n_\tau} \;\leq\; \frac{k_e}{4{,}883{,}809} \;\approx\; 10^{-6} \cdot k_e .
\label{eq:dilution}
\end{equation}
%
Equation~\ref{eq:cost-adjustment} then forces $c_e \lesssim 10^{-6} k_e$
regardless of the value supplied. Every edge becomes almost free, the cost term
in Equation~\ref{eq:pcst-objective} becomes negligible beside the prize term, and
the optimiser has no reason to exclude anything reachable. The prediction is that
the method returns essentially its entire input.

\paragraph{Prediction 2: rank prizes cannot separate identical text.}
Equation~\ref{eq:node-prizes} depends on $\pi(v)$, a strict ordering. When many
nodes share one description their similarities are exactly equal, so $\pi$ is
determined by whatever order the sorting routine happens to produce among ties.
In the measured cohort $4{,}832{,}744$ laboratory tests resolve to $544$ distinct
descriptions, and one patient's $15{,}810$ nodes carry $614$ distinct texts
between them, about $26$ nodes per text. The prediction is that similarity
identifies the
\emph{kind} of entity a question concerns but not \emph{which} instance, and that
two correct implementations of the same algorithm can return different subgraphs
for the same question.

\paragraph{Prediction 3: connectivity is cheap where hubs exist.} The value of
the connectivity constraint depends on paths being expensive. Consider what
happens when a small set of nodes has very high degree, as it does here with
$168$ provenance nodes of average degree $69{,}525$ and six anatomical-site nodes
carrying $3{,}454{,}881$ edges. Almost any pair of nodes is then joined by a path
of length two or three through one of them.
The constraint is then satisfiable at negligible cost, and Steiner nodes stop
being informative. The prediction is that connectivity contributes little until
such nodes are excluded, and that excluding them changes what the subgraph
contains rather than how large it is.

\medskip

None of the three is a defect in G-Retriever, which was designed for and
validated on a different kind of graph. They are the price of the transfer.
Deriving them from Equations~\ref{eq:node-prizes} to~\ref{eq:cost-adjustment},
instead of reporting them as things that happened to be observed, is what makes
them testable predictions rather than post-hoc explanations.
